\documentclass[11pt]{article}

\usepackage[margin=1in]{geometry}
\usepackage[T1]{fontenc}
\usepackage[utf8]{inputenc}
\usepackage{lmodern}
\usepackage{microtype}
\usepackage{amsmath,amssymb,amsthm,mathtools}
\usepackage{booktabs,array}
\usepackage{enumitem}
\usepackage{xcolor}
\usepackage{hyperref}
\usepackage{titlesec}

\hypersetup{
  colorlinks=true,
  linkcolor=blue!50!black,
  citecolor=blue!50!black,
  urlcolor=blue!50!black,
  pdftitle={Two Formalized Partial Results Related to Erd\H{o}s Problem E885},
  pdfauthor={Sam Mausberg}
}

\titleformat{\section}{\Large\bfseries}{\thesection.}{0.6em}{}
\titleformat{\subsection}{\large\bfseries}{\thesubsection.}{0.6em}{}

\newtheoremstyle{plainbreak}
  {\topsep}{\topsep}
  {\itshape}{}
  {\bfseries}{.}
  {\newline}{}
\theoremstyle{plainbreak}
\newtheorem{theorem}{Theorem}[section]
\newtheorem{proposition}[theorem]{Proposition}
\newtheorem{lemma}[theorem]{Lemma}
\newtheorem{corollary}[theorem]{Corollary}

\newtheoremstyle{defbreak}
  {\topsep}{\topsep}
  {}{}
  {\bfseries}{.}
  {\newline}{}
\theoremstyle{defbreak}
\newtheorem{definition}[theorem]{Definition}
\newtheorem{computation}[theorem]{Computation}
\newtheorem{remark}[theorem]{Remark}

\title{Two Formalized Partial Results Related to Erd\H{o}s Problem E885}
\author{Sam Mausberg\\[0.25em]\normalsize Formalized in Lean\,4 with Mathlib}
\date{}

\begin{document}
\maketitle

\begin{center}
\small
\emph{Acknowledgment / Disclosure.}
The source research notes for this writeup were developed with the use of
GPT-5.4 Pro.
The supporting formalization workflow and an initial draft of this note
were prepared with assistance from Aristotle AI
(\url{https://aristotle.harmonic.fun}; GitHub
\texttt{@Aristotle-Harmonic}).
The mathematical claims reported here were then manually curated so that
only statements directly supported by the Lean source remain.
\end{center}

\begin{abstract}
We record two Lean-backed results related to Erd\H{o}s Problem~E885.
First, the explicit triple
\[
  (756000,\;15971200,\;45130176)
\]
has at least five anchored square translates, so the natural universal
anchored cap
\[
  |Y(a,b,c)| \le 4
\]
is false; this does not settle the original $k=5$ problem.
The formalized Lean input here is the fifteen-identity packet
\texttt{new\_3row\_5col\_packet}, and the displayed lower bound is its
immediate finite corollary.
Second, the guidepost seed
\[
  (79200,\;227205,\;1258560)
\]
which is the nonzero part of Choudhry's explicit $(4,4)$ square-sumset
example
\[
  \{0,\;79200,\;227205,\;1258560\}
\]
\cite{Cho25},
is rigid: its common positive factor-difference set is exactly
\[
  \{36,\;468,\;692,\;1028\}.
\]
No claims beyond what is directly supported by the Lean source, together
with immediate finite corollaries of the displayed computations, are made.
\end{abstract}

\section{Introduction}

Erd\H{o}s Problem~E885 asks whether the following holds.
For an integer $n\ge 1$, define the factor-difference set of $n$ by
\[
  D(n) = \{\lvert a-b\rvert : n=ab\}.
\]
Is it true that, for every $k\ge 1$, there exist integers
\[
  N_1 < \cdots < N_k
\]
such that
\[
  \bigl| \bigcap_i D(N_i) \bigr| \ge k?
\]
Background on the problem statement and the previously settled cases
$k=2,3,4$ may be found in
\cite{EP885,ErRo97,Ji99,Br19}.

This note isolates the two strongest standalone results in the present
Lean formalization that are most relevant for a short forum post.
Each referenced Lean theorem or computation below is backed by a
machine-checked Lean\,4 proof, and the cap-failure corollary in
Section~2 is the immediate finite consequence of the displayed five
distinct witnesses.

\section{Anchored 3-shift cap counterexample}

\begin{definition}[Anchored square-translate set]
For positive integers $a,b,c$, define
\[
  Y(a,b,c)
  \;:=\;
  \bigl\{\,z\in\mathbb{Z}_{>0} :
    z^2+a,\ z^2+b,\ z^2+c \text{ are all perfect squares}\,\bigr\}.
\]
\end{definition}

The first result is an explicit packet of five distinct points in one such
set.

\begin{computation}[Five anchored square translates]\label{comp:packet}
For the shift triple
\[
  (a,b,c) \;=\; (756000,\;15971200,\;45130176),
\]
each of the five values
\[
  z \in \{330,\;870,\;2445,\;4155,\;10482\}
\]
lies in $Y(a,b,c)$.
Concretely, the Lean theorem
\texttt{new\_3row\_5col\_packet} proves the fifteen identities
\[
\renewcommand{\arraystretch}{1.15}
\begin{array}{r|ccc}
     & 756000 & 15971200 & 45130176 \\\hline
330  & 930^2   & 4010^2   & 6726^2   \\
870  & 1230^2  & 4090^2   & 6774^2   \\
2445 & 2595^2  & 4685^2   & 7149^2   \\
4155 & 4245^2  & 5765^2   & 7899^2   \\
10482 & 10518^2 & 11218^2 & 12450^2
\end{array}
\]
meaning that for each row value $z$ and each column shift $s$, one has
$z^2+s$ equal to the displayed square.
\end{computation}

\begin{corollary}[Failure of the anchored 3-shift cap]
For the triple $(756000,\;15971200,\;45130176)$ one has
\[
  \{330,\;870,\;2445,\;4155,\;10482\} \subseteq Y(756000,15971200,45130176),
\]
hence
\[
  |Y(756000,15971200,45130176)| \ge 5.
\]
In particular, the universal bound $|Y(a,b,c)|\le 4$ is false.
\end{corollary}
\noindent The formalized input is
\texttt{new\_3row\_5col\_packet} in \texttt{AddendumComputations.lean};
the cardinality conclusion is the immediate finite corollary from the five
displayed distinct witnesses.

\begin{remark}[Why this matters, and what it does not do]
This gives an explicit counterexample to the natural anchored three-shift
bound that had been a plausible intermediate target.
Nearby sumset-in-squares results are adjacent rather than identical to
this anchored positive-shift statement.
Dujella and Elsholtz construct large three-element square-sumsets, and in
their variant where $A$, $B$, and $A+B$ are all squares they reduce to a
zero-shift formulation \cite{DuEl13}.
Elsholtz and Wurzinger later emphasize that $|A|=3$ can support large
$B$, while for $|A|\ge 4$ an absolute bound on $|B|$ is conjecturally
implied by Bombieri--Lang \cite{ElWu24}.
Very recently, Bremner, Elsholtz, and Ulas proved that there are
infinitely many $3$-dimensional Hilbert cubes in the set of squares
\cite{BeElUl26}. This is not the same as the anchored positive-shift
setting considered here, but it is close in spirit.
Choudhry gives explicit $(5,3)$ and $(4,4)$ square-sumset constructions;
displayed examples include
$B=\{0,\;3588000,\;8527200\}$ in the $(5,3)$ case and
$B=\{0,\;79200,\;227205,\;1258560\}$ in the $(4,4)$ case \cite{Cho25}.
The present statement records an explicit positive three-shift anchored
square-translate computation.
It does \emph{not} solve Erd\H{o}s Problem~E885 itself:
the original problem asks for $k$ row values with $k$ common
factor differences, whereas the statement above concerns three fixed
shifts with five anchored square translates.
\end{remark}

\section{Guidepost rigidity}

All named results in this section are formalized in
\texttt{GuidepostRigidity.lean}.

\begin{definition}[Guidepost seed]
For the rigidity discussion we isolate the triple
\[
  (n_1,n_2,n_3) \;=\; (79200,\;227205,\;1258560).
\]
This is the nonzero part of Choudhry's explicit $(4,4)$ square-sumset
example
\[
  \{0,\;79200,\;227205,\;1258560\}
\]
from \cite{Cho25}.
Its row gaps are
\[
  \Delta_1 = n_2-n_1 = 148005,
  \qquad
  \Delta_2 = n_3-n_2 = 1031355.
\]
\end{definition}

The rigidity proof proceeds through an exact finite sieve.

\subsection{Factor-pair images}

Define
\[
  \Sigma(N) \;=\; \{\,d + N/d : d\mid N,\ d < N/d\,\},
  \qquad
  \Delta(N) \;=\; \{\,N/d - d : d\mid N,\ d < N/d\,\}.
\]

\begin{computation}[Guidepost factor-pair images]\label{comp:images}
\begin{align*}
  \Sigma(148005) &= \{774, 794, 838, 922, 954, 1062, 1178, 1382, 1402, 1594,\\
  &\qquad 2214, 2342, 2746, 3334, 3834, 4518, 6458, 9882, 11398,\\
  &\qquad 13466, 16454, 29606, 49338, 148006\},\\[4pt]
  \Delta(1031355) &= \{954, 1062, 1178, 1286, 1402, 2278, 2426, 4582, 4826,\\
  &\qquad 5094, 7866, 8262, 8698, 15802, 22874, 23942, 25114,\\
  &\qquad 26406, 68742, 79322, 114586, 206266, 343782, 1031354\}.
\end{align*}
\end{computation}
\noindent Formalized as \texttt{guidepostSigma\_eq} and
\texttt{guidepostDelta\_eq}.

\subsection{Middle-value sieve}

\begin{lemma}[Guidepost middle-value sieve]\label{lem:middle}
If $d>0$ and
\[
  d \in D(79200)\cap D(227205)\cap D(1258560),
\]
then there exists
\[
  m \in \Sigma(148005)\cap \Delta(1031355)
\]
such that
\[
  m^2 = d^2 + 4\cdot 227205.
\]
\end{lemma}
\noindent Formalized as
\texttt{guidepost\_middle\_value\_of\_common\_secant}.

\begin{computation}[Surviving middle values]\label{comp:middle}
\[
  \Sigma(148005)\cap\Delta(1031355)
  \;=\;
  \{954,\;1062,\;1178,\;1402\}.
\]
\end{computation}
\noindent Formalized as \texttt{guidepostMiddleValues\_eq}.

\subsection{Square test and rigidity}

Each surviving middle value yields a candidate common factor difference via
\[
  d^2 = m^2 - 4\cdot 227205.
\]
The four resulting square tests are:
\begin{align*}
  954^2 - 4\cdot 227205 &= 36^2, \\
  1062^2 - 4\cdot 227205 &= 468^2, \\
  1178^2 - 4\cdot 227205 &= 692^2, \\
  1402^2 - 4\cdot 227205 &= 1028^2.
\end{align*}
These are formalized as \texttt{guidepostSquareTest\_954} through
\texttt{guidepostSquareTest\_1402}.

\begin{theorem}[Guidepost rigidity]\label{thm:rigid}
For every positive integer $d$,
\[
  d \in D(79200)\cap D(227205)\cap D(1258560)
  \quad\Longleftrightarrow\quad
  d \in \{36,\;468,\;692,\;1028\}.
\]
In particular, the guidepost seed has exactly four common positive factor
differences.
\end{theorem}
\noindent The exact biconditional is formalized as
\texttt{guidepost\_positive\_common\_factorDiffSet\_iff}.
The auxiliary Lean theorem \texttt{guidepost\_seed\_rigid} packages the
corresponding finset equality and cardinality statement.

\section{Summary}

This short forum note records two Lean-backed outcomes.
\begin{itemize}[nosep]
  \item The explicit packet in Section~2 implies that the anchored
    three-shift cap is false: the triple
    $(756000,\;15971200,\;45130176)$ has at least five anchored square
    translates.
  \item The guidepost seed
    $(79200,\;227205,\;1258560)$ is rigid: its common positive
    factor-difference set is exactly
    $\{36,\;468,\;692,\;1028\}$.
\end{itemize}
Additional exact identities and explicit $4$-point computations remain in
the repository, but they are omitted here so that the note stays focused on
the strongest forum-facing statements already supported by the Lean source.

\begin{thebibliography}{99}

\bibitem{EP885}
T. Bloom,
\emph{Erd\H{o}s Problem \#885},
\url{https://www.erdosproblems.com/885},
accessed April 18, 2026.

\bibitem{ErRo97}
P. Erd\H{o}s and M. Rosenfeld,
\emph{The factor-difference set of integers},
Acta Arith. \textbf{79} (1997), no.~4, 353--359.

\bibitem{Ji99}
J. Jim\'enez-Urroz,
\emph{A note on a conjecture of Erd\H{o}s and Rosenfeld},
J. Number Theory \textbf{78} (1999), no.~1, 140--143.

\bibitem{Br19}
A. Bremner,
\emph{On a problem of Erd\H{o}s related to common factor differences},
Int. J. Number Theory \textbf{15} (2019), no.~5, 1059--1068.

\bibitem{DuEl13}
A. Dujella and C. Elsholtz,
\emph{Sumsets being squares},
Acta Math. Hungar. \textbf{141} (2013), no.~4, 353--357.
\href{https://doi.org/10.1007/s10474-013-0334-8}{doi:10.1007/s10474-013-0334-8}

\bibitem{ElWu24}
C. Elsholtz and L. Wurzinger,
\emph{Sumsets in the set of squares},
Q. J. Math. \textbf{75} (2024), no.~4, 1243--1254.
\href{https://doi.org/10.1093/qmath/haae044}{doi:10.1093/qmath/haae044}

\bibitem{Cho25}
A. Choudhry,
\emph{Two sets of integers such that all elements of the sumset of the two
sets are perfect squares},
arXiv:2508.07806 [math.NT], 2025.
\href{https://doi.org/10.48550/arXiv.2508.07806}{doi:10.48550/arXiv.2508.07806}

\bibitem{BeElUl26}
A. Bremner, C. Elsholtz, and M. Ulas,
\emph{There are infinitely many Hilbert cubes of dimension 3 in the set of
squares},
arXiv:2604.05459 [math.NT], submitted April 7, 2026.
\href{https://doi.org/10.48550/arXiv.2604.05459}{doi:10.48550/arXiv.2604.05459}

\end{thebibliography}

\end{document}
